\documentclass{hw_template}
\usepackage{xcolor}
\usepackage{graphicx}

% STUDENTS: Please fill these out with your information
\newcommand{\name}{Dylan Andres}
\newcommand{\netid}{dra77}
\newcommand{\collaborators}{}
% / STUDENT

\newcommand{\hwnum}{1}
\newcommand{\duedate}{September 23, 11:59pm ET}
\renewcommand{\title}{Stable Matching, Greedy, Divide and Conquer}
\newcommand{\io}{\textbf{Code input and output format.} }
\newcommand{\submission}{\textbf{Submission.}}
\newcommand{\floor}[1]{\lfloor #1 \rfloor}

% \newtheorem{claim}{Claim}

\begin{document}

%%%%%%%%%%%%%%%%%%%%%%%%%%%%%%%%%%%%%%%%%%%%%%%%%%%%%%%%%%%%%%%%%%%%%%%%%%%%%%%%
% Problem 1: G-S algorithm

\begin{problem}
In this problem, you will compute the maximum sum of non-adjacent elements,
$\operatorname{maxsum}(x)
= \max_{S \in \mathcal{S}} \; \sum_{s \in S} x_s$
where no two elements in \(S\) are adjacent, in two settings: a list and a tree. The definition of
“adjacent” will be specific to each setting. Additionally, for each setting you must first state
the subproblem precisely and then implement the full dynamic programming algorithm.\\


(a) (\textit{2 points})
Given a sequence of integers $x = (x_1, x_2, \ldots, x_n)$,
your goal is to compute the largest sum of non-adjacent elements of \(x\), denoted as
\(\operatorname{maxsum\_list}(x)\).
Adjacent elements of the sequence are next to each other: \(x_i\) and \(x_{i+1}\). \\\\
Identify the dynamic programming subproblem and write its solution here.\\

\textbf{The DP subproblem is: In order to maximize the sum of non-adjacent integers in a list, should we 
take the current integer plus the total sum of the integer two places before it? Or should we
skip the current integer and take the total sum of the integer one place before it?}\\\\

(b) (\textit{2 points})
Given a tree with nodes that contain integers, your goal is to compute the largest sum
of non-adjacent elements. Trees are represented as a set of vertices
$v_i \in V$ and edges $(v_i, v_j) \in E \subseteq V \times V$.
There are no loops and every vertex has at most a single parent.
Adjacent elements of the tree are those that have an edge between nodes, i.e. $(v_i, v_j) \in E$.
For convenience, define the children of a node to be
$c(v_i) = { v_j \mid (v_i, v_j) \in E }$\\\\
Identify the dynamic programming subproblem and write its solution here.\\\\

\textbf{The DP subproblem is: In order to maximum the sum of non-adjacent nodes in a tree, should we
take the current node's value plus the total sum of its sub-trees (excluding its direct child nodes) or should we 
skip the current node's value and take the maximum sum from all of its sub-trees}\\\\

\end{problem}

\newpage
\begin{problem}
Given $n$ points in a 2D plane, your goal is to find the two points that are closest
together according to the Euclidean distance
$d(x, y) = \sqrt{(x_i - y_i)^2 + (x_j - y_j)^2}$.


(d) (\textit{5 points})
Perform an empirical analysis of the three approaches, plotting their runtime
performance for a series of problem sizes. \\\\

\textbf{placeholder}

\end{problem}

\end{document}
